\section{五行状态压缩的可手工核查证书}\label{app:certificate}

第 \ref{sec:five} 节的偶数状态轨道见 \eqref{eq:even-orbits}，
奇数列轨道按以下顺序排列：
\[
 \{1,16\},\{4\},\{7,28\},\{13,22\},\{19,25\},\{31\}.
\]
二十个状态构成从空状态可达的部分。
可达性与封闭性可直接用式 \eqref{eq:Tm} 核对：
从掩码 0 开始逐层列出所有合法输出，得到这些状态，
并且各状态的全部输出仍在该集合中。
行反转交换转移并固定初始状态，故一对反转状态的计数相等。

若 $u_j$ 为偶数列的六维计数列向量，$z_j$ 为下一奇数列的列向量，则
$z_j=Bu_j$，$u_{j+1}=Cz_j$，其中
\begin{equation}\label{eq:BC}
\begin{split}
 B&=\begin{pmatrix}
1&1&1&1&1&0\\
1&2&0&0&0&1\\
1&1&0&0&0&0\\
0&0&0&1&0&0\\
1&0&1&0&0&0\\
1&0&0&0&0&0
\end{pmatrix},\\[4pt]
 C&=\begin{pmatrix}
2&1&2&0&2&1\\
1&1&1&0&0&0\\
1&0&0&0&1&0\\
1&0&0&1&0&0\\
1&0&0&0&0&0\\
0&1&0&0&0&0
\end{pmatrix}.
\end{split}
\end{equation}
为明确系数的含义，若输入轨道为 $O_j$，输出轨道代表为 $q_i$，
则相应单列矩阵项为 $\sum_{p\in O_j}T_5[p,q_i]$。
它与代表的选择无关，但不再要求矩阵对称。
例如 $B_{2,2}=2$：输入掩码 3 和 24 均可合法转移到掩码 4。
这解释了合并后的变量为何取每个状态的计数，而不是轨道计数之和。
直接相乘 $CB$ 得式 \eqref{eq:M5}。

偶数轨道大小向量为 $(1,2,2,2,2,1)$，
奇数轨道大小向量为 $(2,1,2,2,2,1)$。
奇数轨道中行反转固定者是掩码 4 与 31，
偶数轨道中固定者是掩码 0 与 27。
故
\[
 w_{v,o}=(2,1,2,2,2,1)B,\qquad
 w_{r,o}=(0,1,0,0,0,1)B,
\]
与正文所列终端向量一致。

为逐项核对多项式证书，写 $N(t)=n_0+n_1t+n_2t^2+n_3t^3$，其中
\[
\begin{array}{ll}
n_0=(1,0,0,0,0,0)^{\trans},&
n_1=(-7,3,2,1,1,1)^{\trans},\\
n_2=(7,-4,-5,1,0,0)^{\trans},&
n_3=(-1,1,1,0,-1,-1)^{\trans}.
\end{array}
\]
将矩阵 \eqref{eq:M5} 与这四个向量相乘，分别得到
\[
\begin{aligned}
n_1-Mn_0&=-15\evec,&
n_2-Mn_1&=32\evec,\\
n_3-Mn_2&=-15\evec,&
-Mn_3&=\evec.
\end{aligned}
\]
加上常数项 $n_0=\evec$，这些等式正是
$(I-tM)N(t)=(1-15t+32t^2-15t^3+t^4)\evec$ 的全部系数。
因此证书是有限组明确的整数乘法检查，对所有生成函数系数有效。
